\section{Results}
\label{sec:results}
%% ========================================================================

\subsection{Five-Server MCP Architecture}

The system implements five domain-specific MCP servers, each encapsulating a distinct responsibility within the oncology trial workflow. Together, they provide the complete protocol surface required for Physical AI integration with clinical trial systems.

\begin{table}[H]
\centering
\caption{Five MCP servers, their tools, and conformance levels.}
\label{tab:five-servers}
\footnotesize
\begin{tabularx}{\textwidth}{@{}L{3.2cm}C{0.8cm}C{1.0cm}L{6.8cm}@{}}
\toprule
\textbf{Server} & \textbf{Tools} & \textbf{Level} & \textbf{Primary Function} \\
\midrule
\texttt{trialmcp-authz} & 5 & 1 & Deny-by-default RBAC, SHA-256 token lifecycle, policy evaluation \\
\texttt{trialmcp-fhir} & 4 & 2 & FHIR R4 clinical data with HIPAA Safe Harbor de-identification \\
\texttt{trialmcp-dicom} & 4 & 3 & DICOM imaging queries, RECIST~\cite{eisenhauer2009recist} measurements \\
\texttt{trialmcp-ledger} & 5 & 1 & Hash-chained 21 CFR Part 11 audit trail \\
\texttt{trialmcp-provenance} & 5 & 4 & DAG-based data lineage, SHA-256 fingerprinting \\
\midrule
\textbf{Total} & \textbf{23} & & \\
\bottomrule
\end{tabularx}
\end{table}

\subsubsection{Authorization Server}

The authorization server (\texttt{servers/trialmcp\_authz/}) implements deny-by-default role-based access control (RBAC) as defined in ADR-007 (\texttt{docs/adr/007-deny-by-default-rbac.md}). It exposes five tools:

\begin{enumerate}[leftmargin=*]
\item \texttt{authz\_evaluate}: Evaluates whether a given role is permitted to invoke a specific tool.
\item \texttt{authz\_issue\_token}: Issues a scoped, time-limited token with SHA-256 hashing.
\item \texttt{authz\_validate\_token}: Validates a token's authenticity and expiration.
\item \texttt{authz\_revoke\_token}: Revokes a previously issued token.
\item \texttt{authz\_list\_policies}: Lists active policies for a given role or server.
\end{enumerate}

The server is implemented in \texttt{servers/trialmcp\_authz/server.py} (134 lines) with a policy engine (\texttt{policy\_engine.py}) and token store (\texttt{token\_store.py}).

\subsubsection{FHIR Clinical Data Server}

The FHIR server (\texttt{servers/trialmcp\_fhir/}) provides clinical data access with mandatory HIPAA Safe Harbor de-identification. The de-identification pipeline (\texttt{deid\_pipeline.py}) implements HMAC-SHA256 pseudonymization to maintain referential integrity across de-identified records. The four tools are:

\begin{itemize}[leftmargin=*]
\item \texttt{fhir\_read}: Read a single FHIR R4 resource by type and ID.
\item \texttt{fhir\_search}: Search FHIR resources with parameter filtering (limit: 100).
\item \texttt{fhir\_patient\_lookup}: Look up a patient by pseudonymized identifier.
\item \texttt{fhir\_study\_status}: Retrieve the status of a clinical study.
\end{itemize}

\subsubsection{Audit Ledger Server}

The ledger server (\texttt{servers/trialmcp\_ledger/}) implements a hash-chained immutable audit trail satisfying 21 CFR Part 11 requirements~\cite{fda2003part11}. The chain implementation in \texttt{chain.py} uses SHA-256 hashing with canonical JSON serialization. The following code excerpt shows the core hashing mechanism:

\begin{lstlisting}[caption={Hash-chained audit record computation from \texttt{servers/trialmcp\_ledger/chain.py}.},label={lst:chain}]
GENESIS_HASH = "0" * 64

def canonical_json(record: dict[str, Any]) -> str:
    filtered = {k: v for k, v in sorted(record.items())
                if k != "hash"}
    return json.dumps(filtered, sort_keys=True,
                      ensure_ascii=True)

def compute_audit_hash(record, prev_hash: str) -> str:
    canonical = canonical_json(record)
    payload = prev_hash + canonical
    return hashlib.sha256(
        payload.encode("utf-8")).hexdigest()
\end{lstlisting}

Each audit record contains: \texttt{audit\_id}, \texttt{timestamp}, \texttt{server}, \texttt{tool}, \texttt{caller}, \texttt{parameters}, \texttt{result\_summary}, \texttt{previous\_hash}, and \texttt{hash}. The chain is verified by recomputing each hash and comparing it against the stored value.

\subsubsection{Provenance Server}

The provenance server (\texttt{servers/trialmcp\_provenance/}) maintains a directed acyclic graph (DAG) of data lineage records, enabling both forward and backward lineage queries. The DAG implementation in \texttt{dag.py} supports five source types (\texttt{fhir}, \texttt{dicom}, \texttt{computed}, \texttt{external}, \texttt{federated}) and five action types (\texttt{read}, \texttt{transform}, \texttt{aggregate}, \texttt{transfer}, \texttt{derive}).

\subsubsection{DICOM Imaging Server}

The DICOM server (\texttt{servers/trialmcp\_dicom/}) provides imaging data access with role-based permissions and support for RECIST 1.1~\cite{eisenhauer2009recist} measurements critical to oncology trial response evaluation. The server exposes four tools: \texttt{dicom\_query} for searching imaging studies by patient or accession number, \texttt{dicom\_retrieve\_pointer} for obtaining secure references to image data without transferring pixel data, \texttt{dicom\_study\_metadata} for retrieving study-level DICOM metadata, and \texttt{dicom\_recist\_measurements} for querying tumor measurement data used in treatment response assessment.

The DICOM adapter (\texttt{dicom\_adapter.py}) supports integration with industry-standard PACS systems through the 9 DICOM integration adapters in \texttt{integrations/dicom/}, including dcm4chee, Orthanc, and DICOMweb connectors.

\subsubsection{Server Architecture Diagram}

The following text diagram shows the complete five-server deployment topology with data flow paths:

\begin{figure}[H]
\begin{verbatim}
  NATIONAL MCP FIVE-SERVER DEPLOYMENT
  +---------------------------------------------------+
  |                  SITE DEPLOYMENT                  |
  |                                                   |
  |    +-----------+  +-----------+  +-----------+    |
  |    | trialmcp  |  | trialmcp  |  | trialmcp  |    |
  |    | authz     |  | fhir      |  | dicom     |    |
  |    | (5 tools) |  | (4 tools) |  | (4 tools) |    |
  |    | RBAC,     |  | FHIR R4,  |  | DICOMweb, |    |
  |    | Tokens    |  | HIPAA     |  | RECIST    |    |
  |    +-----+-----+  +-----+-----+  +-----+-----+    |
  |          |              |              |          |
  |          +------+-------+------+-------+          |
  |                 |              |                  |
  |          +------v------+ +-----v-------+          |
  |          | trialmcp    | | trialmcp    |          |
  |          | ledger      | | provenance  |          |
  |          | (5 tools)   | | (5 tools)   |          |
  |          | SHA-256     | | DAG, W3C    |          |
  |          | 21 CFR 11   | | PROV        |          |
  |          +-------------+ +-------------+          |
  |                                                   |
  |  23 tools total across 5 servers                  |
  +---------------------------------------------------+
\end{verbatim}
\end{figure}

\subsubsection{Complete Tool Inventory}

Table~\ref{tab:tool-inventory} provides the complete inventory of all 23 tools with their input parameters and response types:

\begin{table}[H]
\centering
\caption{Complete 23-tool inventory across all five MCP servers.}
\label{tab:tool-inventory}
\footnotesize
\begin{tabularx}{\textwidth}{@{}L{4.2cm}L{3.5cm}L{4.1cm}@{}}
\toprule
\textbf{Tool} & \textbf{Key Parameters} & \textbf{Response Type} \\
\midrule
\texttt{authz\_evaluate} & role, tool & ALLOW/DENY decision \\
\texttt{authz\_issue\_token} & role, expires\_in & Token with SHA-256 hash \\
\texttt{authz\_validate\_token} & token\_hash & Valid/invalid status \\
\texttt{authz\_revoke\_token} & token\_hash & Revocation confirmation \\
\texttt{authz\_list\_policies} & role, server & Policy list \\
\midrule
\texttt{fhir\_read} & resource\_type, id & De-identified FHIR resource \\
\texttt{fhir\_search} & resource\_type, params & De-identified bundle \\
\texttt{fhir\_patient\_lookup} & pseudonym & De-identified patient \\
\texttt{fhir\_study\_status} & study\_id & Study status object \\
\midrule
\texttt{dicom\_query} & patient\_id, modality & Study result list \\
\texttt{dicom\_retrieve\_pointer} & study\_uid & Secure image pointer \\
\texttt{dicom\_study\_metadata} & study\_uid & DICOM metadata \\
\texttt{dicom\_recist\_measurements} & study\_uid & RECIST measurements \\
\midrule
\texttt{ledger\_append} & server, tool, caller & Audit record with hash \\
\texttt{ledger\_verify} & (none) & Chain validity status \\
\texttt{ledger\_query} & server, tool, caller & Filtered audit records \\
\texttt{ledger\_replay} & start\_index, count & Record sequence \\
\texttt{ledger\_export} & format & Full chain export \\
\midrule
\texttt{provenance\_record} & source\_id, action, actor & Provenance record \\
\texttt{provenance\_query\_forward} & record\_id, depth & Downstream lineage \\
\texttt{provenance\_query\_backward} & record\_id, depth & Upstream lineage \\
\texttt{provenance\_fingerprint} & data & SHA-256 fingerprint \\
\texttt{provenance\_lineage} & source\_id & Full lineage graph \\
\bottomrule
\end{tabularx}
\end{table}

\subsection{Conformance Levels and Profiles}

The standard defines five hierarchical conformance levels, each building upon the previous:

\begin{table}[H]
\centering
\caption{Five hierarchical conformance levels with required servers and use cases.}
\label{tab:conformance-levels}
\small
\begin{tabularx}{\textwidth}{@{}C{0.7cm}L{2.2cm}L{3.8cm}L{5cm}@{}}
\toprule
\textbf{Level} & \textbf{Name} & \textbf{Required Servers} & \textbf{Use Case} \\
\midrule
1 & Core & authz + ledger & Basic authorization and audit \\
2 & Clinical Read & Level 1 + fhir & FHIR clinical data access \\
3 & Imaging & Level 2 + dicom & DICOM imaging queries \\
4 & Federated Site & Level 3 + provenance & Multi-site data lineage \\
5 & Robot Procedure & All 5 servers & Full autonomous clinical workflow \\
\bottomrule
\end{tabularx}
\end{table}

Eight conformance profiles are defined in \texttt{profiles/}, including five hierarchical profiles and three regulatory overlays (\texttt{country-us-fda.md}, \texttt{state-us-ca.md}, \texttt{state-us-ny.md}).

\subsection{Safety Modules for Physical AI}

The system provides eight safety modules in \texttt{safety/} specifically designed for robotic systems operating in clinical environments:

\begin{table}[H]
\centering
\caption{Safety modules for Physical AI systems.}
\label{tab:safety-modules}
\footnotesize
\begin{tabularx}{\textwidth}{@{}L{3.8cm}L{8.2cm}@{}}
\toprule
\textbf{Module} & \textbf{Purpose} \\
\midrule
\texttt{estop.py} & Emergency stop signal propagation, abort workflow, post-abort evidence capture, recovery with re-authorization \\
\texttt{procedure\_state.py} & Procedure state machine enforcement (IDLE, PREPARING, ACTIVE, PAUSED, COMPLETING, ABORTED) \\
\texttt{robot\_registry.py} & Robot registration, capability validation, firmware version tracking \\
\texttt{task\_validator.py} & Task order validation against safety constraints and robot capabilities \\
\texttt{gate\_service.py} & Safety gate evaluation before procedure transitions \\
\texttt{approval\_checkpoint.py} & Multi-party approval checkpoints (clinician + system + regulatory) \\
\texttt{site\_verifier.py} & Site capability verification against conformance requirements \\
\bottomrule
\end{tabularx}
\end{table}

The emergency stop controller in \texttt{safety/estop.py} implements a four-state lifecycle (IDLE, TRIGGERED, ACKNOWLEDGED, RECOVERING) with signal propagation to all registered servers. This ensures that any authorized party can halt a robotic procedure immediately, with full state preservation and evidence capture.

\subsection{Test Coverage and Validation}

The repository contains 668 test functions across two test suites:

\begin{table}[H]
\centering
\caption{Test coverage summary across unit and conformance suites.}
\label{tab:test-coverage}
\small
\begin{tabularx}{\textwidth}{@{}L{3cm}R{1.5cm}L{7cm}@{}}
\toprule
\textbf{Category} & \textbf{Tests} & \textbf{Validates} \\
\midrule
Unit (integrations) & 205 & All 34 integration adapters \\
Unit (safety) & 88 & All 8 safety modules \\
Unit (core server) & 32 & 5 MCP server implementations \\
Unit (infrastructure) & 12 & Conformance harness, schema validator \\
Positive conformance & 69 & Correct behavior at Levels 1--3 \\
Adversarial & 57 & AuthZ bypass, chain tampering, PHI leakage, replay attacks \\
Blackbox & 63 & Per-server external conformance \\
Security & 39 & SSRF prevention, token lifecycle, chain integrity \\
Interoperability & 32 & Cross-server trace, schema validation \\
Negative & 26 & Invalid input rejection, unauthorized access \\
Integration & 13 & Multi-server integration scenarios \\
\midrule
\textbf{Total} & \textbf{668} & \\
\bottomrule
\end{tabularx}
\end{table}

\subsection{Integration Adapters}

The repository implements 34 integration adapter modules across six categories in \texttt{integrations/}:

\begin{table}[H]
\centering
\caption{Integration adapter categories and module counts.}
\label{tab:integration-adapters}
\small
\begin{tabularx}{\textwidth}{@{}L{2cm}R{1.2cm}L{8.5cm}@{}}
\toprule
\textbf{Category} & \textbf{Count} & \textbf{Key Adapters} \\
\midrule
FHIR & 9 & HAPI FHIR, SMART on FHIR, de-identification, terminology, patient filter \\
DICOM & 9 & dcm4chee, Orthanc, DICOMweb, RECIST, modality filter \\
Clinical & 3 & Scheduling, e-consent, provenance export \\
Federation & 4 & Coordinator, policy enforcement, secure aggregation \\
Identity & 5 & OIDC, mTLS, KMS, policy engine \\
Privacy & 4 & De-identification pipeline, privacy budget, data residency \\
\midrule
\textbf{Total} & \textbf{34} & \\
\bottomrule
\end{tabularx}
\end{table}

\subsection{Repository Scale}

Table~\ref{tab:repo-scale} summarizes the quantitative scope of the national repository at version 1.2.0:

\begin{table}[H]
\centering
\caption{Repository quantitative summary (v1.1.0).}
\label{tab:repo-scale}
\small
\begin{tabular}{@{}lr|lr@{}}
\toprule
\textbf{Category} & \textbf{Count} & \textbf{Category} & \textbf{Count} \\
\midrule
MCP servers & 5 & Safety modules & 8 \\
Tools (total) & 23 & Benchmark modules & 5 \\
JSON schemas & 13 & Certification tools & 4 \\
Conformance profiles & 8 & Deployment modes & 6 \\
Test functions & 668 & Interop scenarios & 8 \\
Integration adapters & 34 & Actor personas & 6 \\
Specification docs & 9 & Stakeholder guides & 5 \\
ADRs & 7 & Operational docs & 5 \\
Total files & 381 & Total directories & 88 \\
Lines of code (approx.) & 69,800 & Regulatory overlays & 7 \\
\bottomrule
\end{tabular}
\end{table}

\subsection{Deployment Infrastructure}

The system supports six deployment configurations, enabling evaluation at any scale:

\begin{enumerate}[leftmargin=*]
\item \arraybackslash \raggedright \textbf{Individual Docker containers:} Five Dockerfiles in \texttt{deploy/docker/} (e.g., \texttt{Dockerfile.authz}, \texttt{Dockerfile.fhir}).
\item \textbf{All-in-one Docker:} \texttt{Dockerfile.allinone} for single-container evaluation.
\item \textbf{Docker Compose (single-site):} \texttt{deploy/docker-compose.yml} for local deployment.
\item \textbf{Docker Compose (multi-site):} \texttt{deploy/docker-compose.multi-site.yml} for federated testing.
\item \textbf{Kubernetes manifests:} 7 manifests in \texttt{deploy/kubernetes/} for production clusters.
\item \textbf{Helm chart:} \texttt{deploy/helm/trialmcp/} with \texttt{Chart.yaml} and \texttt{values.yaml}.
\end{enumerate} \raggedright

Server configuration is managed through YAML files in \texttt{deploy/config/} (one per server), enabling site-specific customization of storage backends, authentication providers, and privacy settings.

\subsection{SDK and Developer Tooling}

The repository provides dual-language SDKs and developer tools to facilitate adoption:

\begin{itemize}[leftmargin=*]
\item \textbf{Python SDK} (\texttt{sdk/python/trialmcp\_client/}): 8 client modules with 4 middleware components (audit, auth, circuit breaker, retry). Installable via \texttt{pip install -e .} with \texttt{pyproject.toml}.

\item \textbf{TypeScript SDK} (\texttt{sdk/typescript/src/}): 8 client modules with matching middleware stack. Configured via \texttt{package.json} and \texttt{tsconfig.json}.

\item \textbf{CLI} (\texttt{tools/cli/trialmcp\_cli.py}): Commands for \texttt{init}, \texttt{scaffold}, \texttt{validate}, \texttt{certify}, \texttt{schema diff}, and \texttt{config generate}.

\item \textbf{Code generation} (\texttt{tools/codegen/}): Generators for Python models, TypeScript interfaces, and OpenAPI specifications from JSON schemas.
\end{itemize}

Both SDKs include role-scoped example scripts for all six actor types: robot agent, trial coordinator, data monitor, auditor, sponsor, and CRO.

\subsection{Schema and Specification Infrastructure}

The repository defines 13 machine-readable JSON schemas in \texttt{schemas/} using JSON Schema Draft 2020-12. These schemas serve as the normative contract between MCP server implementations and clients, enabling automated validation and code generation:

\begin{table}[H]
\centering
\caption{JSON schemas organized by domain category.}
\label{tab:schemas}
\scriptsize
\begin{tabularx}{\textwidth}{@{}L{4.5cm}L{1.5cm}L{5.5cm}@{}}
\toprule
\textbf{Schema File} & \textbf{Category} & \textbf{Validates} \\
\midrule
\texttt{audit-record.schema.json} & Audit & Audit trail entries \\
\texttt{authz-decision.schema.json} & AuthZ & RBAC evaluation results \\
\texttt{capability-descriptor.schema.json} & Infra & Server capability declarations \\
\texttt{consent-status.schema.json} & Privacy & Patient consent state \\
\texttt{dicom-query.schema.json} & Imaging & DICOM query parameters \\
\texttt{error-response.schema.json} & Infra & Standardized error envelopes \\
\texttt{fhir-read.schema.json} & Clinical & FHIR resource read responses \\
\texttt{fhir-search.schema.json} & Clinical & FHIR search bundle responses \\
\texttt{health-status.schema.json} & Infra & Server health check responses \\
\texttt{provenance-record.schema.json} & Prov & DAG lineage records \\
\texttt{robot-capability-profile} & Robotics & Robot system capabilities \\
\texttt{site-capability-profile} & Infra & Site-level capabilities \\
\texttt{task-order.schema.json} & Robotics & Procedure task orders \\
\bottomrule
\end{tabularx}
\end{table}

Nine normative specification documents in \texttt{spec/} define the complete protocol surface, including \texttt{core.md} (protocol scope, design principles, 5 conformance levels), \texttt{actor-model.md} (6 actors with permission matrices), \texttt{tool-contracts.md} (23 tool signatures with input/output contracts), \texttt{security.md} (RBAC, mTLS, OIDC, SSRF prevention), and \texttt{privacy.md} (HIPAA Safe Harbor, de-identification, consent, data residency).

\subsection{Emergency Stop Implementation}

The emergency stop controller is a critical safety component. The following code excerpt from \texttt{safety/estop.py} shows the four-state lifecycle and signal propagation:

\begin{lstlisting}[caption={Emergency stop lifecycle from \texttt{safety/estop.py}.},label={lst:estop}]
class EStopStatus(Enum):
    IDLE = "IDLE"
    TRIGGERED = "TRIGGERED"
    ACKNOWLEDGED = "ACKNOWLEDGED"
    RECOVERING = "RECOVERING"

class EStopController:
    def trigger_estop(self, procedure_id, reason,
                      triggered_by, current_state=None):
        if self._status == EStopStatus.TRIGGERED:
            raise RuntimeError(
                "E-stop already active")
        signal = EStopSignal(
            signal_id=str(uuid.uuid4()),
            procedure_id=procedure_id,
            reason=reason,
            triggered_by=triggered_by,
            affected_servers=list(
                self._active_servers),
            preserved_state=current_state or {})
        self._active_signal = signal
        self._status = EStopStatus.TRIGGERED
        return signal
\end{lstlisting}

The e-stop lifecycle ensures: (1) only one e-stop can be active at a time, (2) all registered servers receive the stop signal, (3) procedure state is preserved for post-incident analysis, and (4) recovery requires explicit re-authorization. These properties are validated by 88 unit tests in \texttt{tests/test\_safety.py}.

\subsection{Interoperability Testbed}

The interoperability testbed in \texttt{interop-testbed/} provides a multi-site simulation environment with 8 scenarios and 6 actor personas. Each scenario validates a specific cross-site interaction pattern:

\begin{enumerate}[leftmargin=*]
\item \texttt{audit\_replay.py}: Verifies audit trail replay across sites
\item \texttt{cross\_site\_provenance.py}: Validates cross-site provenance chain integrity
\item \texttt{partial\_outage.py}: Tests graceful degradation during partial failures
\item \texttt{robot\_workflow.py}: End-to-end robotic procedure workflow across servers
\item \texttt{schema\_drift.py}: Tests schema version compatibility under evolution
\item \texttt{site\_onboarding.py}: Validates new site onboarding workflow
\item \texttt{state\_overlay.py}: Tests state-level regulatory overlay application
\item \texttt{token\_exchange.py}: Validates cross-site token exchange protocols
\end{enumerate}

The testbed uses Docker Compose (\texttt{interop-testbed/docker-compose.yml}) with mock services for EHR (\texttt{mock\_ehr.py}), PACS (\texttt{mock\_pacs.py}), and identity (\texttt{mock\_identity.py}) to simulate realistic clinical environments.

%% ========================================================================
